\section{Maturity of the Project}\label{sec:maturity}

\subsection{Development stage}\label{sec:stage}

CarbonPass is an early prototype purpose-built for this competition; it has not
been commercialized, and the modification-rule provisions therefore do not
apply. At submission the following is demonstrable:

\begin{itemize}[leftmargin=1.3em,itemsep=2pt]
\item \textbf{Module 1 (Data-Pack Engine)} runs end-to-end on a three-firm mock
corpus built from real Taiwanese Taipower bill and Ministry of Finance
e-invoice formats, producing a per-CN-code data pack in the European Commission
template with per-line uncertainty.
\item \textbf{Module 2 (Grid-Aware Scheduler)} runs on a synthetic
fastener-plant load profile against the \emph{live} open 10-minute grid feed,
computing hourly carbon intensity and an optimized shift plan\cite{taipowerdata}.
\item \textbf{Module 3 (Inclusion Interface)} is an operational LINE
conversational front end in Mandarin/Taiwanese.
\end{itemize}

The proof of concept requires no hardware, no paid API, and no partnership to
run --- it is built entirely on open data (Section~\ref{sec:opendata}), which
makes its feasibility independently checkable by the jury.

\subsection{Project adjustments}\label{sec:adjustments}

Not applicable: the project is original to this competition and has not been
sold or awarded elsewhere. No 50\% modification was required.

\subsection{Demonstration}\label{sec:demo}

A demonstration video accompanies the online application, showing the full
path from photographed documents to a generated data pack, the buyer-cost delta
screen, and a grid-aware schedule computed from the live open grid feed.
\emph{[Video link entered in application form field 3.3.]}

\subsection{Existing validation fields}\label{sec:valfields}

Validation is planned within the competition's own mentorship machinery and
involves the public sector. During the mentorship period, outreach is directed
to a TIFI member factory in the Gangshan cluster for a real pilot, with the
Metal Industries Research \& Development Centre (MIRDC) and Kaohsiung
Environmental Protection Bureau's voluntary-reduction advisory team as domain
and public-sector validation partners. \emph{These are validation targets being
approached, not yet secured commitments}, and are stated as such.

\subsection{Existing validation results and verification plan}\label{sec:valresults}

At submission the quantified results are those produced on the mock corpus and
synthetic-plus-live-feed pipeline (Section~\ref{sec:stage}); the field-measured
results are generated during mentorship and presented at the final review. The
plan maps directly onto the official calendar\cite{handbook}:

\begin{table}[htbp]
\caption{Implementation and verification milestones mapped to the official schedule}
\label{tab:plan}
\small
\begin{tabularx}{\textwidth}{@{}p{0.27\textwidth}X@{}}
\toprule
\textbf{Window (Handbook)} & \textbf{Milestone} \\
\midrule
By 31 Jul 2026 (submission) & Proof of concept v0 + demonstration video + this proposal; validation outreach opened to TIFI, MIRDC, and Kaohsiung EPB \\
6--16 Aug (preliminary review) & Proposal and working prototype under evaluation \\
Mid-Sep -- mid-Oct (mentorship) & One pilot factory in Gangshan: real bills in, real data pack out; pack structure reviewed with an MIRDC engineer or accredited verifier; one flexible load shifted and measured on-meter; weekly progress log \\
Late Oct (final review) & Live demonstration: data pack generated on stage from raw documents; live grid-aware schedule; pilot before/after ledger; market outlook updated with the September European Parliament vote\cite{akin} \\
\bottomrule
\end{tabularx}
\end{table}

\begin{formfield}{Boundaries stated up front}
CarbonPass \emph{prepares} verification; it does not replace the accredited
verifier --- at least 80\% of reported emissions must be third-party
verified\cite{omm}. Preferential carbon-fee rates belong to large regulated
emitters, not SMEs; the SME's incentives are customer retention, cost
pass-through, energy savings, and readiness\cite{moenvfee}. The 2028 scope
extension is pending legislation and is treated as a tracked roadmap event, not
a booked benefit\cite{akin}.
\end{formfield}
